
  \section{Related work}
  \label{sec:related}

  Two lines of work bear on the research question: LLMs placed on top of an existing
  simulation or DT, and LLMs that call tools in sequence, search or plan.

  \subsection{LLMs over simulations and DTs}
  \label{sec:related-dt}

  Kleiman et al.\ build a framework on the pairing of an LLM with a simulation, placing an agent between the user
  and an existing simulation \cite{kleiman2025simulationagent}. Such a layer adds \emph{access},
  turning the user's question into a correctly configured run, and \emph{interpretation}, turning
  the result into an explanation the user can act on, without the model needing to know the
  process itself. Two of the framework's design choices carry over to this project
  \cite{kleiman2025simulationagent}: the simulation is called but never modified, and its raw
  output is reduced to compact summaries before the model sees it, because a full time series
  exceeds what a model can usefully consume.

  Other work gives the model more to decide. Ndum et al.\ add an LLM to the DT
  of an advanced nuclear reactor. It calls the twin's tools and draws on domain knowledge to
  suggest operating steps, while the operator keeps the decision \cite{ndum2026aroma}. Xia et al.\
  let LLM agents run a twin's simulation repeatedly to find parameter settings that
  meet an objective, and show this in a case study, leaving systematic evaluation for later
  \cite{xia2024llmsimulation}. This is the closest precedent for the inverse questions here: the
  model itself searches over a simulation it cannot change.

  Mylonas et al.\ go the other way, in an agentic twin for a power-grid control room
  \cite{mylonas2026governance}. Their model chooses tools from an approved list and sets their
  arguments, but code outside its control writes the numbers into each answer from the tools'
  results and makes any action with side effects wait for an operator's approval. These rules
  held in every run of a 118-task benchmark, and without that code only 39.2\% of answers had
  numbers the tools supported. Where an answer's numbers come from is thus
  measured there, but enforced in code. In this project it rests on the skills' instructions,
  which the model can depart from, and that is why part (a) of the research question measures it.

  \subsection{Tool use, search and planning}
  \label{sec:related-tools}

  Tool use has moved from single calls to sequences. ReAct interleaves reasoning with actions,
  choosing each call from the result of the last \cite{yao2023react}, and program-aided models
  write code for an interpreter so that arithmetic is done there rather than in generated text
  \cite{gao2023pal}. A model can also act as the search itself, proposing new candidates from the
  scores of earlier ones \cite{yang2024opro}. An inverse question needs all three: a sequence of
  calls, numbers kept out of generated text, and a search the model steers.

  Whether the model should run that search itself is less settled. On problems from the
  automated-planning community, Valmeekam et al.\ find that LLMs generate poor plans,
  the strongest models included \cite{valmeekam2023planbench}. Guan et al.\ therefore have the
  model write a formal model of the domain and leave the search to a sound planner, with the model
  translating between the planner and the user \cite{guan2023worldmodels}. The heater-boost
  advisory also keeps the search out of the model (\Cref{sec:boost}): its search for the lowest
  boost is coded in the tool layer, and a shorter opening is suggested but not searched for. For an inverse question with
  no tool behind it, the choice is between coding such a search in advance and leaving it to the
  model, and the research question compares the two.

  \subsection{Position of this work}
  \label{sec:related-position}

  None of these studies compares a model's own search over a twin's simulation with a search coded
  over the same tools, in accuracy or in cost, and the one that measures where an answer's numbers
  come from does so behind a layer the model cannot bypass \cite{mylonas2026governance}. This project does both for a layer
  made only of skills, which instruct the model but do not constrain it. It measures how close the
  model's searched answers come to a coded search and whether it says which parts the tools did not
  produce, and what leaving the search to the model costs and saves compared with coding it
  (\Cref{sec:rq}).
